\documentclass[a4paper,11pt]{article}

\usepackage{fontspec}
\usepackage{etoolbox}
\usepackage[margin=25mm]{geometry}
\usepackage{xcolor}
\usepackage{hyperref}
\usepackage{microtype}
\usepackage{needspace}

\setmainfont{Latin Modern Roman}
\setsansfont{Latin Modern Sans}
\setmonofont[
  BoldFont = LMMono10-Regular,
  Scale = 0.9
]{Latin Modern Mono}

\emergencystretch=2em
\AtBeginEnvironment{verbatim}{\linespread{1.0}\selectfont}

\hypersetup{
  pdftitle  = {OLLM -- OSG LaTeX Lecture Maker},
  pdfauthor = {OSG},
  colorlinks,
  linkcolor = blue!50!black,
  urlcolor  = blue!50!black
}

\title{OLLM\\OSG LaTeX Lecture Maker}
\author{OSG}
\date{Entwurfsstand}

\begin{document}

\maketitle

\begin{abstract}
  OLLM ist das Build-Frontend des \texttt{osglecture}-Bundles. Es wählt
  Dokument- und Sprachvarianten aus und übergibt den einzelnen LaTeX-Build an
  \texttt{latexmk}. Dieses Handbuch entsteht parallel zur Neuimplementierung.
\end{abstract}

\tableofcontents

\section{Status}

Das ausgelieferte Programm \texttt{ollm} besitzt bereits den neuen portablen
Kommandozeilen-Launcher. Kompatible Builds übergibt er zunächst an die
erhaltene \texttt{latexmk}-Konfiguration \texttt{ollm-legacy.rc} in Version
0.11.1. Ihre häufig verwendete Schnittstelle bildet die Kompatibilitätsbasis
für die Überarbeitung.

Die Zielarchitektur mit ihren Dateiverträgen und Rationales steht in
\texttt{DESIGN.md}. Noch nicht implementierte Schnittstellen werden in diesem
Handbuch nicht als verfügbar beschrieben.

\section{Aufgabenteilung}

\begin{description}
  \item[OLLM] wählt und orchestriert konkrete Builds.
  \item[\texttt{latexmk}] führt LuaLaTeX und benötigte Hilfsprogramme bis zum
    stabilen Ergebnis aus.
  \item[\texttt{osglecture}] interpretiert den Buildauftrag und bestimmt
    Dokumentsemantik, Serienstruktur und Referenzverhalten.
  \item[Deploymentwerkzeug] soll später geprüfte Artefakte veröffentlichen.
\end{description}

\section{Legacy-Aufruf}

OLLM wird direkt aufgerufen. Der Launcher startet für kompatible Builds
\texttt{latexmk} mit der Legacy-RC-Datei. Beispiele der bisherigen Oberfläche
sind:

\begin{verbatim}
ollm slides
ollm handout
ollm script
ollm +lang=en script
ollm standalone
\end{verbatim}

Die genaue neue CLI wird erst dokumentiert, wenn ihr Verhalten durch Tests
festgeschrieben ist.

\section{Neue Grundfunktionen}

Die normalisierte Auswahl kann bereits ohne LaTeX-Build ausgegeben werden:

\begin{verbatim}
ollm build script --language=en --dry-run
ollm build script --language=en --dry-run --format=json
\end{verbatim}

\texttt{ollm doctor} prüft, ob Perl, \texttt{latexmk} und LuaLaTeX auffindbar
sind. \texttt{ollm --version} gibt die Version des neuen Launchers aus.

\section{Projektmanifest}

Das Projektmanifest heißt \texttt{ollmconfig.toml} und markiert die
Projektwurzel. OLLM sucht ausgehend vom Arbeitsverzeichnis aufwärts danach.
Alternativ werden Manifest oder Projektwurzel ausdrücklich ausgewählt:

\begin{verbatim}
ollm build script --config=/path/to/ollmconfig.toml --dry-run
ollm build script --project-root=/path/to/project --dry-run
\end{verbatim}

Existieren \texttt{ollmconfig.toml} und das alte \texttt{ollmconfig.pl} im
selben Verzeichnis, bricht OLLM ab. Damit hängt die gewählte
Konfigurationssemantik nicht von einer stillen Prioritätsregel ab.

Ein minimales Serienmanifest enthält:

{\small
\begin{verbatim}
schema = 1
profile = "OSG lecture/1"

[project]
id = "bs"
title = "Betriebssysteme"

[languages]
available = ["de", "en"]
default = "de"

[languages.map]
de = "ngerman"
en = "british"

[targets.slides]
languages = ["de"]

[targets.script]
languages = ["de", "en"]

[security]
shell_escape = "restricted"
\end{verbatim}
}

\subsection*{Profile und Targets}

Profile und Dokumenttargets sind versionierte Definitionen. OLLM löst sie
zunächst aus dem mitgelieferten Bundle auf. Zusätzliche Suchpfade werden in
der optionalen Datei \texttt{.ollmconfig.local.toml} eingetragen:

\begin{verbatim}
schema = 1

[definitions]
paths = ["configuration", "/opt/osglecture/definitions"]
\end{verbatim}

Relative Pfade beziehen sich auf die Projektwurzel. Unbekannte Schlüssel,
Profile und Targets sind Fehler. Soweit die TOML-Quelle einen Schlüssel
enthält, nennt die Diagnose auch dessen Zeile. Technische Schemawerte sind
unabhängig von der Bundleversion; bei einem Konflikt nennt OLLM gefundenes und
unterstütztes Schema ausdrücklich.

\section{Buildauftrag}

Für einen konkreten Serienbuild normalisiert OLLM die Verzeichnisidentität und
erzeugt einen an den Jobnamen gebundenen Auftrag:

\begin{verbatim}
<jobname>.osgbuild.tex
\end{verbatim}

Die Datei enthält unter anderem Dokumenttyp, Sprache, physische Einheit,
Projektprofil und eine SHA-256-Konfigurationssignatur. Sie wird zunächst
temporär im Zielverzeichnis geschrieben und nur bei verändertem Inhalt
ersetzt. Dadurch lösen unveränderte Aufträge keinen unnötigen LaTeX-Neubau aus.

Das Paket \texttt{osglecture-config} liest den Auftrag und lehnt ihn ab, wenn
sein Schema nicht unterstützt wird oder die enthaltene Job-ID nicht mit dem
aktuellen TeX-Jobnamen übereinstimmt. Die Klasse verwendet Dokumenttyp und
Sprache aus diesem expliziten Auftrag; nur der Legacyzweig liest sie weiterhin
aus dem Jobnamen.

\section{TOML-Parser}

OLLM liefert den reinen Leseweg von \texttt{TOML::Tiny} 0.22 mit. Nutzer
müssen daher kein CPAN-Modul nachinstallieren. \texttt{ollm doctor} berichtet
Name, Version und Herkunft des aktiven Parsers.

Diese Bündelung ist für TeX Live wichtig: Drittmodule werden für
Anwendungsskripte nicht allgemein garantiert, und insbesondere das
mitgelieferte Windows-Perl ist bewusst minimal. OLLM verwendet nur
\texttt{Parser}, \texttt{Tokenizer} und \texttt{Grammar}; diese benötigen
ausschließlich Perl-Core-Module.

Für Entwicklung und Vergleich kann die vollständige Upstream-Distribution
optional installiert werden:

\begin{verbatim}
cpanm TOML::Tiny
cpan TOML::Tiny
\end{verbatim}

OLLM verwendet weiterhin die mitgelieferte und getestete Fassung. Eine
CPAN-Installation ist daher weder Voraussetzung noch Reparaturschritt.

\Needspace{10\baselineskip}
\section{Entwicklungs-Build}

Das OLLM-Modul ist in den Bundle-Build mit \texttt{l3build} integriert:

\begin{verbatim}
l3build check
l3build doc
l3build install
\end{verbatim}

\texttt{check} prüft Perl-Syntax, CLI, Manifestvalidierung, lokale
Definitionspfade sowie die Auflösung des bundleweiten Minimalbeispiels.

\section{Weiterführende Dokumente}

\begin{description}
  \item[\texttt{README.md}] kurzer Einstieg und Entwicklungsbefehle;
  \item[\texttt{DEPENDENCIES.md}] Laufzeitabhängigkeiten und
    Installationshinweise;
  \item[\texttt{DESIGN.md}] Architektur, Invarianten, Rationales und offene
    Entscheidungen;
  \item[\texttt{THIRD\_PARTY.md}] Herkunft und Lizenz mitgelieferter
    Komponenten;
  \item[dieses Handbuch] nutzerorientierte Beschreibung implementierter
    Funktionen.
\end{description}

\end{document}
